\subsection{The group of similarities.}

Let \(\Phi\) be a Hermitian form on E. An automorphism \(u\) of the A-module E is called a similarity (relative to \(\Phi\) ) if there exists an invertible element \(\alpha\) of A such that one has

\[
\Phi (u (x), u (y)) = \alpha \Phi (x, y)\tag{7}
\]

for all x, y in E. The similarities form a group Γ.

When \(\Phi\) takes values that are regular elements of A (for example when A is a field and \(\Phi \neq 0\) ), the element \(\alpha\) of A satisfying (7) is uniquely determined by \(u\) ; it is called the multiplier of the similarity \(u\) . Changing \(x\) to \(\lambda x\) in (7), we then see that \(\alpha\) belongs to the center of A; exchanging \(x\) and \(y\) in (7), we see moreover that \(\overline{\alpha} = \alpha\) . If, for \(u \in \Gamma\) , we denote by \(\alpha(u)\) the multiplier of \(u\) , the map \(u \to \alpha(u)\) is a homomorphism of \(\Gamma\) into the multiplicative group of invertible elements of the center of A. The kernel of this homomorphism is the unitary group associated with \(\Phi\) , which is therefore a normal subgroup of \(\Gamma\) . Let \(\beta\) be an invertible element of the center of A, \(v\) the homothety of ratio \(\beta\) , and \(w\) a unitary automorphism of E; then \(vw = wv\) is a similarity of E, and its multiplier is \(\beta\overline{\beta}\) . Conversely, let \(u\) be a similarity whose multiplier is of the form \(\beta\overline{\beta}\) ( \(\beta\) denoting an invertible element of the center of A); then \(uv^{-1}\) is a unitary automorphism \(w\) , hence \(u\) is of the form \(vw\) .

Suppose now that A is a field, E is a finite-dimensional vector space, and that \(\Phi\) is nondegenerate. For every similarity \(u\) with multiplier \(\alpha\) we have

\[
\Phi (x, \alpha y) = \alpha \Phi (x, y) = \Phi (u (x), u (y)) = \Phi (x, u ^ {*} (u (y))),
\]

therefore \(u^*u\) is the homothety with ratio \(\alpha\) . If A is commutative, and if \(n\) denotes the dimension of E, it follows from this and from formula (50) of § 1, no. 10 that one has

\[
(\det u) \overline {{(\det u)}} = \alpha^ {n}.\tag{8}
\]

Let us then distinguish two cases:

1°) The integer n is odd, i.e., \(n = 2q + 1\) . Then, setting \(\rho = \alpha^{-q}(\det u)\) , one has \(\alpha = (\det u) \overline{(\det u)} \alpha^{-2q} = \rho \bar{\rho}\) . Thus u is the product of the homothety with ratio \(\rho\) and a unitary automorphism.

2°) The integer n is even, say n = 2q. Then, setting \(\rho = \alpha^{-q}(\det u)\) , one has \(\rho\bar{\rho} = 1\) . In particular, when J is the identity, one has \((\det u)^{2} = \alpha(u)^{2q}\) ; the similarities u such that det \(u = \alpha(u)^{q}\) (resp. det \(u = -\alpha(u)^{q}\) ) are called direct (resp. inverse); the direct similarities form a normal subgroup of index 2 of \(\Gamma\) ;

the homotheties with ratio ≠ 0 are direct similarities; the same is true of orthogonal transformations with determinant 1 (no. 2); orthogonal transformations with determinant -1 are inverse similarities.

The preceding definitions and results remain valid for ε-hermitian forms (§ 3, no. 1), and in particular for alternating forms.

Let A be a commutative field and Q a quadratic form \(\neq 0\) on E. One calls a similarity (relative to Q) any automorphism \(u\) of E such that there exists a nonzero element \(\alpha\) of A (called the multiplier of \(u\) ) for which \(Q(u(x)) = \alpha Q(x)\) for every \(x \in E\) . It is clear that \(u\) is then a similarity with multiplier \(\alpha\) relative to the bilinear form associated with Q; the converse is true when the characteristic of A is \(\neq 2\) .
% semantic-fragments: legacy-input-seam
\relax{}
